\section{Deployment and operational administration}

\subsection{Why deployment remains a human-executed task}

Most C-LARA-2 work takes place inside the repository. Codex can inspect
existing material, write and revise code, add tests, maintain documentation
and update the project-memory records described in the preceding section.
Deployment crosses a different boundary. It requires actions in AWS, SSH
sessions, process-management tools, browser or cloud-provider interfaces, and
production-server environments. For safety and practical reasons, human
collaborators execute and authorise those external actions; Codex does not hold
unrestricted production access.

This does not mean that humans must already possess the relevant system
administration knowledge. Codex can inspect repository and deployment context,
explain the purpose of an operation, formulate concrete diagnostic or repair
steps, interpret returned results and preserve the resulting procedure. The
human role is to perform, supervise and accept the external operation. The AI
role is broad operational guidance in a domain that is only partly accessible
from the repository itself.

\subsection{One-off provisioning and recurrent administration}

Two kinds of operational documentation support the AWS deployment. The
one-off task of provisioning, configuring and validating the server is
recorded in the deployment-and-migration roadmap. Like other roadmaps, it
captures the intended outcome, design choices, migration and validation steps,
recovery considerations and lessons learned while the work was carried out.
It is not expected to be followed repeatedly from beginning to end.

Recurrent tasks are documented separately in the server-administration
how-to file. The most important of these is installing a new version of the
platform and checking that the running service is healthy afterwards. The
separation is useful: a historical roadmap records why the deployment was
constructed in a particular way, while a concise runbook gives human
collaborators a repeatable procedure for operating it. Both documents are
maintained by Codex in response to human needs and operational evidence.

\subsection{The read-only Assistant sandbox}

One important case, discussed in more detail in Section~9, is the
project-understanding Assistant. The Assistant must be able to inspect the
current repository and answer questions about C-LARA-2 with evidence, but it
must not receive unrestricted authority to alter the repository or the
production server. Deploying this capability therefore required a read-only,
restricted execution environment.

The most complicated operational work concerned configuring that environment
using Bubblewrap\footnote{Bubblewrap is a Linux sandboxing tool that can run a
process in a restricted filesystem and operating-system environment rather
than giving it the ordinary permissions of the surrounding server.} and
AppArmor\footnote{AppArmor is a Linux security framework that applies profiles
restricting what a program may read, write, execute, or otherwise access.}.
In broad terms, these mechanisms were needed because repository inspection by
a coding agent is useful only if it can be separated from broader powers to
modify files, access unrelated server material, execute unsafe commands or
affect the running service.

The human collaborators did not merely lack the commands needed to configure
these mechanisms; they had not previously encountered Bubblewrap or AppArmor
or understood their role in the deployment. Codex investigated failure modes,
interpreted diagnostics, identified relevant constraints, proposed and revised
procedures, and recorded the resulting knowledge in version-controlled
documentation. Humans performed the external operations and checked the
results.

This episode illustrates an important qualification to the phrase ``human
supervision.'' Humans retained credentials, production access and final
acceptance, but they could not independently reconstruct the specialist
systems-administration reasoning behind every action. They therefore placed
bounded trust in Codex's technical judgement, using restricted permissions,
diagnostics, reversibility, documentation and post-change checks to make that
trust inspectable. The result was not autonomous AI control of the AWS server;
it was a form of human--AI cooperation in which human authority was retained
but technical competence was substantially delegated.

\subsection{Lessons for delegated responsibility}

The deployment episode may foreshadow a broader pattern in AI-assisted work.
As coding assistants become capable of reasoning across codebases,
infrastructure, diagnostics and operational documentation, human collaborators
may increasingly retain formal authority while relying on AI systems for
technical judgement that they cannot independently reproduce. The important
research question is not whether such reliance occurs, but how it can be made
appropriately bounded, transparent, auditable and reversible.

The deployment work provides a concrete example of responsibility expanding
through demonstrated competence. Codex was initially used to guide humans
through tasks they could not otherwise have planned or explained. As it
successfully analysed failures, revised procedures and preserved operational
knowledge, it became reasonable to delegate a larger part of the sysadmin role
to it, while humans retained production access and final acceptance. This does
not establish that every strategic, cultural or ethical decision should be
transferred to an AI. It does show that responsibility need not be fixed.
Within a bounded domain, delegation can increase when competence, constraints,
evidence, logging and human oversight make the arrangement justified. C-LARA-2
therefore offers one early case in which this arrangement made a difficult
project goal achievable for a team that otherwise lacked the required
specialist expertise.
